\chapter{ヘッジファンドとは}
\label{ch:hf}

\section{定義}

ヘッジファンド（hedge fund）には、法律で決まった定義はない。ふつうは、次の特徴をいくつか持つ運用の器（ファンド）をそう呼ぶ。

\begin{defbox}[ヘッジファンドの典型的な特徴]
\begin{enumerate}
\item \textbf{私募}：広く一般に売らず、富裕層・年金・大学基金・保険会社などの限られた投資家から資金を集める。
\item \textbf{手段の自由}：空売り、借入れ（レバレッジ）、先物・オプション・スワップなどのデリバティブを使える。投資信託（ミューチュアルファンド）は、これらを法律や約款で強く制限されている。
\item \textbf{絶対リターンの目標}：市場の指数に勝つこと（相対リターン）ではなく、市場が上がっても下がってもプラスの収益を上げること（絶対リターン）を目標に掲げることが多い。
\item \textbf{成功報酬}：運用の利益の一部（伝統的には 20\%）を運用者が受け取る。
\item \textbf{解約の制限}：毎日は解約できず、月・四半期ごとに、事前の通知をして解約する。最初の 1 年は解約できない（ロックアップ）こともある。
\end{enumerate}
\end{defbox}

「ヘッジ」は、もともと「市場全体の上げ下げのリスクを、空売りで打ち消す」という意味である。最初のヘッジファンドとされる A.~W.~Jones の運用（1949 年）は、割安と考える株を買い、割高と考える株を空売りして、市場の動きに左右されにくくしたうえで、借入れで規模を大きくした。いまの「ヘッジファンド」は、市場のリスクを打ち消さない戦略（たとえば通貨や金利の方向に大きく賭けるグローバル・マクロ）も含む、広い言葉になっている。

\section{歴史の流れ}

\begin{table}[H]
\centering\small
\caption{ヘッジファンドの歴史の主な出来事}
\label{tab:history}
\begin{tabular}{@{}lP{0.74\linewidth}@{}}
\toprule
年 & 出来事 \\
\midrule
1949 & A.~W.~Jones が、ロング・ショートとレバレッジを組み合わせた投資組合を始める。利益の 20\% を報酬とした。 \\
1966 & 経済誌 \textit{Fortune} の記事が Jones の成績を紹介し、同じ形のファンドが増える。 \\
1970 年代 & 1973〜74 年の下げ相場で多くのファンドが損をして閉じる。Bridgewater（1975）などが設立される。 \\
1980 年代 & グローバル・マクロの時代。Soros の Quantum Fund、Robertson の Tiger など。Renaissance Technologies（1982）、D.~E.~Shaw（1988）など、数学・計算機を使う運用会社が生まれる。 \\
1992 & Soros のファンドが英ポンドの売りで大きな利益を上げる（欧州の為替相場メカニズムからの英国の離脱）。 \\
1998 & LTCM（Long-Term Capital Management）が、ロシアの債務不履行をきっかけに巨額の損失を出し、ニューヨーク連銀の仲介で銀行団が救済する。 \\
2000 年代 & 年金・大学基金の資金が流入し、業界が急拡大する。AQR（1998）、Two Sigma（2001）などのクオンツ運用会社が成長する。 \\
2007 & 8 月に株式のクオンツ戦略が数日で急落する（quant quake）。 \\
2008 & 金融危機。Lehman Brothers の破綻でプライムブローカーに預けた資産が凍結され、多くのファンドが解約の制限（ゲート）をかける。 \\
2010 年代 & 手数料の引き下げの圧力、ファクター投資（スマートベータ）の普及、複数の運用チームを抱えるマルチマネージャーの拡大。 \\
2020 年代 & 2020 年 3 月の急落と回復、2021 年の個人投資家による空売りの踏み上げ（GameStop）、代替データと機械学習の普及、LLM の登場。 \\
\bottomrule
\end{tabular}
\end{table}

業界の規模は、データ会社や業界団体の推計で、2020 年代半ばに運用資産 4 兆ドル前後とされる（推計の方法によって幅がある）。ただし、その資産の多くは大手の数十社に集中している。

\section{仕組み}

\subsection{法的な形}

米国のヘッジファンドは、たいてい次の形をとる（\figref{fig:structure}）。

\begin{itemize}
\item \textbf{リミテッド・パートナーシップ}：運用会社（または関連会社）が無限責任のジェネラル・パートナー（GP）、投資家が有限責任のリミテッド・パートナー（LP）になる。
\item \textbf{投資会社法の適用の除外}：米国の 1940 年投資会社法の規制を避けるため、投資家の数を 100 人以下に限るか（3(c)(1)）、一定の資産を持つ「適格購入者」に限る（3(c)(7)）。広告をせずに私募で売る。
\item \textbf{マスター・フィーダー}：米国の課税される投資家のための米国内のファンド（フィーダー）と、米国外の投資家や非課税の投資家のための国外（ケイマン諸島など）のファンド（フィーダー）が、同じマスター・ファンドに資金を入れ、マスターがまとめて売買する。
\item \textbf{運用会社の登録}：2010 年の Dodd--Frank 法の後、一定の規模を超える運用会社は SEC に投資顧問として登録し、ファンドの規模やレバレッジを報告する（Form PF）。先物を主に使う運用会社（CTA）は CFTC の規制も受ける。
\end{itemize}

日本では、金融商品取引法の適格機関投資家等特例業務などの枠組みで、私募のファンドを運用する。

\begin{figure}[H]
\centering
\begin{tikzpicture}[node distance=7mm and 9mm]
\node[box=c4] (inv1) {米国の課税\\投資家};
\node[box=c4,right=of inv1] (inv2) {米国外・非課税\\の投資家};
\node[box=c1,below=of inv1] (f1) {米国内\\フィーダー（LP）};
\node[box=c1,below=of inv2] (f2) {国外フィーダー\\（ケイマン等）};
\node[box=c7,below=11mm of $(f1)!0.5!(f2)$,text width=30mm] (m) {マスター・ファンド\\（売買をまとめる）};
\node[box=c3,left=24mm of m] (gp) {運用会社\\（GP・投資顧問）};
\node[box=c2,right=24mm of m] (pb) {プライム\\ブローカー};
\node[box=ink2,below=8mm of m,text width=30mm] (mk) {市場（株・債券・\\先物・デリバティブ）};
\draw[arr] (inv1)--(f1);
\draw[arr] (inv2)--(f2);
\draw[arr] (f1)--(m);
\draw[arr] (f2)--(m);
\draw[arr] (gp)--node[lab,above]{運用の指示}(m);
\draw[arr,<->] (m)--node[lab,above]{資金・株の貸借}(pb);
\draw[arr,<->] (m)--(mk);
\end{tikzpicture}
\caption{マスター・フィーダーの構造。ほかに、基準価額を計算する管理会社（アドミニストレーター）、監査法人、法律事務所がかかわる。}
\label{fig:structure}
\end{figure}

\subsection{プライムブローカー}

ヘッジファンドの売買を支えるのが、投資銀行のプライムブローカー部門である。

\begin{itemize}
\item \textbf{資金の貸付け}：買う株を担保に資金を貸す（信用買い）。
\item \textbf{株の貸付け}：空売りのための株を探して貸す（ロケート）。借りる株には貸株料がかかり、借り手が多い株ほど高い。貸し手が株を返せと言えば（リコール）、空売りを閉じなければならない。
\item \textbf{保管・決済}：売買の決済と、資産の保管。
\item \textbf{証拠金の計算}：ポジションのリスクから必要な担保（証拠金）を計算し、損失が出れば追加の担保を求める（マージンコール）。
\end{itemize}

2008 年の Lehman Brothers の破綻では、その英国法人にプライムブローカーとして預けていた資産が長く凍結された。この後、多くのファンドが複数のプライムブローカーを使うようになった。

\subsection{解約の条件}

ヘッジファンドの資産には、すぐに売れないもの（小型株、社債、私募の資産など）が含まれる。投資家が一斉に解約すると、安値で投げ売りしなければならない。それを防ぐため、次のような条件をつける。

\begin{itemize}
\item \textbf{解約の頻度と通知の期間}：例えば「四半期末に、45 日前までに通知」。
\item \textbf{ロックアップ}：最初の 1 年などは解約できない。解約に手数料をとる「ソフト・ロックアップ」もある。
\item \textbf{ゲート}：1 回の解約で出せる額を、ファンド全体の一定の割合までに制限する。
\item \textbf{サイドポケット}：売れない資産を別の勘定に分け、売れるまで解約の対象から外す。
\end{itemize}

資産の流動性と、投資家の解約の条件が合っていないと、危機のときに投げ売りが起きる。これは第 II 部で扱う失敗の多くに共通する構造である。

\section{手数料}
\label{sec:fees}

\subsection{2 と 20}

伝統的な手数料は「2 と 20」と呼ばれる。

\begin{itemize}
\item \textbf{管理報酬}（management fee）：運用資産の年 2\%。成績にかかわらずかかる。
\item \textbf{成功報酬}（performance fee、incentive fee）：利益の 20\%。
\end{itemize}

実際の平均はこれより低く、2020 年代には管理報酬が 1\% 台半ば、成功報酬が 10\% 台後半と報じられることが多い。一方で、評判の高い一部のファンドは、もっと高い報酬をとる。複数の運用チームを抱える大手のマルチマネージャーでは、運用チームの報酬・データ・システムなどの費用をそのまま投資家に回す「パススルー」の形が広がり、投資家の負担はむしろ大きくなっている。

\subsection{ハイウォーターマークとハードル}

成功報酬には、次の決まりがつくことが多い。

\begin{itemize}
\item \textbf{ハイウォーターマーク}：基準価額が過去の最高値を超えた分にだけ成功報酬をとる。損をした翌年に取り返しただけでは、成功報酬はかからない。
\item \textbf{ハードル}：短期金利などの一定の収益を超えた分にだけ成功報酬をとる。
\end{itemize}

\begin{exbox}[手数料の計算]
投資家が 100 を預け、手数料の前の収益が年 10\% だったとする（短期金利 0、ハードルなし、ハイウォーターマークの下ではない）。
\begin{enumerate}
\item 管理報酬：$100 \times 2\% = 2$。
\item 成功報酬：手数料の前の利益 10 から管理報酬を引いた 8 の 20\%、つまり $1.6$。
\item 投資家の取り分：$10 - 2 - 1.6 = 6.4$、つまり 6.4\%。
\end{enumerate}
手数料の前の収益の 36\% が運用者に渡る。手数料の前の収益が 5\% なら、$5 - 2 - 0.6 = 2.4$ で、半分以上が運用者に渡る。手数料の前の収益が $-10\%$ なら、管理報酬で $-12\%$ になり、成功報酬はゼロである。
\end{exbox}

成功報酬は、運用者にとって「ファンドの利益に対するコールオプション」のような形をしている。損をしても運用者は損を負わず、利益が出れば一部を受け取る。そのため、リスクを取りすぎる誘因があると指摘されてきた。ハイウォーターマークの下に大きく沈んだファンドは、成功報酬を当分受け取れないので、閉じて別のファンドを立ち上げる誘因も生まれる。

\section{運用会社の型}

\subsection{シングル・マネージャーとマルチマネージャー}

\begin{itemize}
\item \textbf{シングル・マネージャー}：一人（または一つのチーム）の運用者の判断や、一つの戦略の体系で運用する。
\item \textbf{マルチマネージャー}（プラットフォーム型、ポッド型）：数十〜数百の小さな運用チーム（ポッド）に資金を割り当て、会社の中央でリスクを管理する。各チームには損失の上限があり、数\% の損失で割り当てを減らされ、さらに損をするとチームが解散させられる、と報じられる。市場全体やスタイル・ファクター（\chapref{ch:factor}）のリスクは打ち消すよう求められ、チームの「固有の」収益（アルファ）だけを積み上げる設計になっている。
\item \textbf{ファンド・オブ・ファンズ}：複数のヘッジファンドに投資するファンド。分散と選別を売りにするが、手数料が二重にかかる。
\end{itemize}

\subsection{裁量とシステマティック}

\begin{itemize}
\item \textbf{裁量（discretionary）}：運用者が、企業の調査、経済の見通し、経験に基づいて売買を判断する。
\item \textbf{システマティック（systematic、クオンツ）}：売買のルールを、データと統計・最適化のモデルで決め、計算機で実行する。人は、モデルを作り、検証し、監視する。
\end{itemize}

実際には、その中間（クオンツのスクリーニングで候補を絞り、人が判断する「クオンタメンタル」など）も多い。この報告書は、主にシステマティックな方法を扱う。

\subsection{ほかの運用との違い}

\begin{table}[H]
\centering\small
\caption{ヘッジファンドと、ほかの運用の違い（典型的な場合）}
\label{tab:types}
\begin{tabular}{@{}P{0.17\linewidth}P{0.18\linewidth}P{0.18\linewidth}P{0.18\linewidth}P{0.18\linewidth}@{}}
\toprule
 & 投資信託 & ヘッジファンド & 自己勘定の取引会社 & CTA（先物の運用） \\
\midrule
資金 & 一般の投資家 & 富裕層・機関投資家 & 会社自身 & 機関投資家・一般 \\
目標 & 指数に勝つ・連動 & 絶対リターン & 会社の利益 & 絶対リターン \\
空売り・借入れ & 強く制限 & 自由 & 自由 & 先物で自由 \\
報酬 & 資産の 0.1〜1\% 台 & 管理報酬＋成功報酬 & 社員の歩合 & 管理報酬＋成功報酬 \\
解約 & 毎日 & 月・四半期 & なし & 月・毎日 \\
典型的な戦略 & ロングのみの株・債券 & 多様 & 高頻度・マーケットメイク & トレンドフォロー \\
\bottomrule
\end{tabular}
\end{table}

\section{収益はどこから来るか}

ヘッジファンドの収益を、次の三つに分けて考えると見通しがよい（\chapref{ch:perf}、\chapref{ch:factor} で式にする）。

\begin{enumerate}
\item \textbf{市場のベータ}：株式市場全体が上がれば上がる部分。指数のファンドで安く手に入る。
\item \textbf{オルタナティブ・ベータ}（スタイル・プレミアム、リスク・プレミアム）：割安株、モメンタム、キャリー、トレンドフォロー、ボラティリティの売りなど、ルールで再現できる戦略の収益。昔はヘッジファンドの「腕」とされたが、研究が進み、いまは低い手数料の商品でも手に入る。
\item \textbf{アルファ}：上の二つで説明できない部分。情報、分析、速さ、執行の上手さ、ほかの人が取れないリスク（流動性の低い資産を持つ、混雑した取引の反対に立つ）から来る。
\end{enumerate}

研究で繰り返し問われてきたのは、「ヘッジファンドの手数料の後の収益のうち、本当のアルファはどれだけか」である。第 II 部（第 \ref{ch:hfperf} 章）で見るように、多くの研究は、平均的なファンドの収益の大部分が 1 と 2 で説明でき、手数料の後のアルファは小さいか、時代とともに小さくなったと報告している。

\begin{projbox}
LLM-BLM のポートフォリオは、ロングのみで、時価総額加重からの乖離（TE）に上限をつけた。これはヘッジファンドというより、指数に対して運用する伝統的な株式の運用会社（ロングのみのアクティブ運用）の設計に近い。空売りが使えないので、悪いと読んだ銘柄からは「持たない」分しか収益を取れない。この違いは、\chapref{ch:pc} と第 \ref{ch:implications} 章で再び扱う。
\end{projbox}
